% Part: proof-theory
% Chapter: proof-search
% Section: search-algorithm

\documentclass[../../../include/open-logic-section]{subfiles}

\begin{document}

\olfileid{pt}{ps}{sal}

\olsection{\Log{G3c} لئی بھال دا الگورتھم}

اسیں \Log{G3c} لئی بھال دا اک الگورتھم دَسّنے آں۔ ایہو اکّو ممکن الگورتھم نہیں، تے سبھ توں کارگر وی نہیں۔ پر ایہ درست اے۔ جے سیکوئنٹ $\Gamma \Sequent \Delta$ دا !!a{proof} اے، تاں ایہ آخر رُک جاندا اے تے !!a{proof} لبھ لیندا اے۔ ایہ خطا وی نہیں کھاندا۔ جے !!a{proof} لبھے بغیر رُک جاوے، یا کدے رُکے ای نہ، تاں $\Gamma
\Sequent \Delta$ شروع توں ای معتبر نہیں۔

الگورتھم دی اک ضروری خاصیت ایہ اے کہ ایہ پکا کردا اے کہ ہر !!{formula} جیہڑا $\Gamma \Sequent \Delta$ وچ، یا بھال دوران بنے کسے سیکوئنٹ وچ اے، اوہدی «تقلیل» ہووے۔ یعنی اوہنوں پچھاں ول لائے گئے استنباطی قاعدے دا مُکھ فارمولا بنایا جاوے، تے ایہ جِنّی واری لوڑ ہووے اُنی واری ہووے۔ ایس خاصیت نوں \emph{نرپکھتا} آکھدے آں۔

الگورتھم پڑاؤواں وچ کم کردا اے۔ ہر پڑاؤ دا حاصل سیکوئنٹاں دا اک درخت اے، تے اگلے پڑاؤ دا حاصل، اوہناں سارے سبھ توں اُتلے سیکوئنٹاں وچ !!a{formula} دی تقلیل کر کے بناندا اے، جیہڑے پہلاں ای مسلّمے نہیں۔

پڑاؤ $0$ اُتے $\Gamma \Sequent \Delta$ توں شروع کردے آں۔ نرپکھتا پکی کرن لئی اسیں !!{formula}s، $ \Gamma$ تے $ \Delta$ وچلے، نوں انڈیکس دے طور اُتے عدد دیندے آں، ایویں کہ وکھ وکھ !!{formula}s نوں وکھ وکھ عدد ملن۔ مثال لئی !!{formula}s نوں کھبّے توں سجّے ول گنّن بارے سوچو۔ عدد اُتلے نشان دے طور اُتے لکھدے آں۔ مثال لئی جے اسیں !!a{proof} لبھ رہے آں $!A, !B \Sequent !C$ دا، تاں پڑاؤ~$0$ دا حاصل $!A^1, !B^2 \Sequent !C^3$ اے۔ ایہو جیہے انڈیکس والے سیکوئنٹ وچ اُتلے نشان دے طور اُتے آؤن والا سبھ توں وڈّا~$n$ اوہدا \emph{سبھ توں وڈّا انڈیکس} اے (مثال وچ~$3$)۔

جے سیکوئنٹ وچ مسوّر نیں، تاں ایہ وی پتا ہونا چاہیدا اے کہ سجّے پاسے $ \lexists[x][!A(x)]$ یا کھبّے پاسے $ \lforall[x][!A(x)]$ دی تقلیل ویلے کِہڑے ٹرم ورتنے نیں۔ سانوں صرف !!{constant}s $ \Obj c_0$، $ \Obj c_1$، $ \Obj c_2$، \dots، توں بنے ٹرم چاہیدے نیں، جیہڑے~$\Gamma$ تے $ \Delta$ وچ آؤن والیاں کسے وی فنکشن نشانیاں نال بنن۔ اسیں منّنے آں کہ ایہناں سارے ٹرماں دا سیٹ کسے مقرر گنتی وار لڑی وچ دِتّا اے، تاں جو اسیں مثلاً «پہلے !!{constant} جیہڑا $ \Gamma \Sequent \Delta$ وچ نہیں آؤندا» دی گل کر سکیے۔ الگورتھم دے بنائے درخت وچ ہر انڈیکس والے سیکوئنٹ نال !!{constant}s دا اک سیٹ جُڑیا اے۔ پڑاؤ~$0$ اُتے سیکوئنٹ نوں دِتّا !!{constant}s دا سیٹ اوہ سارے !!{constant}s رکھدا اے جیہڑے اوہدے وچ آؤندے نیں (جے کوئی ہون)، تے $\{\Obj c_0\}$ اے جے کوئی !!{constant}s نہیں۔

منّ لو کہ الگورتھم نے پڑاؤ~$n$ اُتے انڈیکس والے سیکوئنٹاں دا درخت بنایا اے، نال !!{constant}s دے جُڑے سیٹ۔ پڑاؤ~$n+1$ اُتے اسیں ہر سبھ توں اُتلے سیکوئنٹ لئی ہِٹھ لکھیا کردے آں۔

جے !!a{formula}~$!A$ سیکوئنٹ دے کھبّے تے سجّے دونے پاسے آؤندا اے، یا کھبّے پاسے~$\lfalse$ اے، تاں اسیں کجھ نہیں کردے۔ ایہو جیہے سیکوئنٹاں دے !!{proof}s~\Log{G3c} وچ نیں، تے ایس خاص سبھ توں اُتلے سیکوئنٹ لئی ہور کم نہیں رہندا۔

جے سیکوئنٹ وچ کوئی غیر جوہری !!{formula} نہیں، تاں الگورتھم وچّے ای روک دیو۔ $\Gamma \Sequent \Delta$ دا کوئی ثبوت نہیں۔

نئیں تاں غیر جوہری !!{formula}~$!A$ چُنو جیہدا انڈیکس~$i$ سبھ توں چھوٹا اے۔ تاں ایہ سبھ توں اُتلا سیکوئنٹ یا $!A^i, \Pi
\Sequent \Lambda$ اے یا $ \Pi \Sequent \Lambda, !A^i$۔ $k$ نوں سیکوئنٹ دا سبھ توں وڈّا انڈیکس تے $C$ نوں اوہنوں دِتّا !!{constant}s دا سیٹ منّ لو۔

اسیں $!A^i$ دی «تقلیل» کردے آں، یعنی استنباطی قاعدے دے حساب نال نویں سبھ توں اُتلے سیکوئنٹ بناؤنے آں۔ ایہ قاعدہ !!{main operator}، جیہڑا~$!A$ دا اے، تے ایہ گل کہ $!A^i$ کھبّے پاسے اے یا سجّے، طے کردے نیں۔

\begin{enumerate}
  \item $!A \ident !B \land !C$ تے کھبّے پاسے آؤندا اے۔ $(!B \land !C)^i, \Pi \Sequent \Lambda$ توں اُتّے اک نواں سبھ توں اُتلا سیکوئنٹ شامل کرو:
  \[
  \Axiom$!B^k, !C^{k+1}, \Pi \fCenter \Lambda$
  \RightLabel{\LeftR{\land}}
  \UnaryInf$(!B \land !C)^i, \Pi \fCenter \Lambda$
  \DisplayProof
  \]
  نویں سیکوئنٹ نوں دِتّا !!{constant}s دا سیٹ~$C$ اے۔
  \item $!A \ident !B \land !C$ تے سجّے پاسے آؤندا اے۔ $\Pi \Sequent \Lambda, (!B \land !C)^i$ توں اُتّے دو نویں سبھ توں اُتلے سیکوئنٹ شامل کرو:
  \[
  \Axiom$\Pi \fCenter \Lambda, !B^k$
  \Axiom$\Pi \fCenter \Lambda, !C^k$
  \RightLabel{\RightR{\land}}
  \BinaryInf$\Pi \fCenter \Lambda, (!B \land !C)^i$
  \DisplayProof
  \]
  نویں سیکوئنٹ نوں دِتّا !!{constant}s دا سیٹ~$C$ اے۔

  \item $!A \ident \lnot !B$، $!A \ident !B \lor !C$ تے $!A \ident !B \lif !C$ والیاں صورتاں ایہناں ورگیاں نیں تے مشقاں لئی چھڈّیاں گئیاں نیں۔

  \item $!A \ident \lexists[x][!B(x)]$ تے کھبّے پاسے آؤندا اے۔ $\lexists[x][!B(x)]^i, \Pi \Sequent \Lambda$ توں اُتّے اک نواں سبھ توں اُتلا سیکوئنٹ شامل کرو:
  \[
  \Axiom$!B(c)^k, \Pi \fCenter \Lambda$
  \RightLabel{\LeftR{\lexists}}
  \UnaryInf$\lexists[x][!B(x)]^i, \Pi \fCenter \Lambda$
  \DisplayProof
  \]
  جتھے $c$ پہلا !!{constant} اے جیہڑا~$C$ وچ نہیں۔ نویں سیکوئنٹ نوں دِتّا !!{constant}s دا سیٹ $C \cup \{c\}$ اے۔
  \item $!A \ident \lexists[x][!B(x)]$ تے سجّے پاسے آؤندا اے۔ $\Pi \Sequent \Lambda, \lexists[x][!B(x)]$ توں اُتّے اک نواں سبھ توں اُتلا سیکوئنٹ شامل کرو:
  \[
  \Axiom$\Pi \fCenter \Lambda, \lexists[x][!B(x)]^k, !B(t)^{k+1}$
  \RightLabel{\LeftR{\lexists}}
  \UnaryInf$\Pi \fCenter \Lambda, \lexists[x][!B(x)]^i$
  \DisplayProof
  \]
  جتھے $t$ پہلا اوہ ٹرم اے جیہدے وچ صرف اوہ !!{constant}s ہون جیہڑے~$C$ وچ نیں، تے جیہڑا ایس سیکوئنٹ توں تھلّے سجّے پاسے $ \lexists[x][!B(x)]$ دی تقلیل لئی پہلاں نہ ورتیا ہووے (یا $\Obj c_0$، جے ایہو جیہے سارے ٹرم ورتے جا چکے نیں)۔ نویں سیکوئنٹ نوں دِتّا !!{constant}s دا سیٹ $C$ اے۔
  \item $!A \ident \lforall[x][!B(x)]$ والیاں صورتاں ایہناں ورگیاں نیں تے مشقاں لئی چھڈّیاں گئیاں نیں۔
\end{enumerate}

جے پڑاؤ $n+1$ اُتے اسیں کسے وی سبھ توں اُتلے سیکوئنٹ نال کوئی نواں سیکوئنٹ نہیں جوڑیا، تاں الگورتھم کامیابی نال ختم ہو جاندا اے۔ ایس صورت وچ اسیں !!a{proof} لبھ لیا اے~$\Gamma \Sequent \Delta$ دا، جیہڑا پڑاؤ $n+1$ دے درخت توں سارے انڈیکس ہٹا کے ملدا اے۔ سبھ توں اُتلے سارے سیکوئنٹ یا $!A, \Pi \Sequent \Lambda, !A$ دی شکل دے نیں یا $\lfalse, \Pi
\Sequent \Lambda$ دی۔ سارے استنباط~\Log{G3c} دے درست استنباط نیں۔ قضیاتی استنباطاں لئی ایہ صاف اے۔ ویکھو کہ ہر قدم اُتے سیٹ~$C$، یعنی سیکوئنٹ نوں دِتّا !!{constant}s دا سیٹ، اوہ سارے !!{constant}s رکھدا اے جیہڑے سیکوئنٹ وچ آؤندے نیں۔ ایس لئی \LeftR{\lexists} تے \RightR{\lforall} نال تقلیل وچ خاص متغیر دی شرط پوری ہوندی اے۔ خاص متغیر~$c$ اک !!a{constant} سی جیہڑا نتیجے نوں دِتّے سیٹ~$C$ وچ نہیں سی، سو $c$ نتیجے وچ وی نہیں آؤندا سی۔ سانوں ایہ ملدا اے:

\begin{prop}
\Log{G3c} لئی ثبوت لبھّن دا الگورتھم درست اے۔ جے ایہ کامیابی نال ختم ہووے، تاں شروع دے سیکوئنٹ~$\Gamma \Sequent \Delta$ دا !!a{proof} اے۔
\end{prop}

\end{document}
